% PAGES: 319-321

see (10.40). Note that $D^2F_0(x_0)$ is an $n \times n$ symmetric matrix.

We restrict our attention to the case when the matrix $D^2F_0(x_0)$ is either
symmetric positive definite, or symmetric negative definite, which suffices for the
purpose of this book. The general theory of Lagrange multipliers can be found in
section 9.1 from Stefanica \cite{36}; see Theorem 9.2 therein.

\begin{theorem}
Assume that the function $g(x)$ satisfies condition (10.66). Let $x_0 \in U \subset
\R^n$ and $\lambda_0 \in \R^m$ such that the point $(x_0,\lambda_0)$ is a critical
point for the Lagrangian function $F(x,\lambda) = f(x) + \lambda^tg(x)$. Let
$F_0(x) = f(x) + \lambda_0^tg(x)$, and let $D^2F_0(x_0)$ be the Hessian of $F_0$
evaluated at the point $x_0$.

If the matrix $D^2F_0(x_0)$ is positive definite, i.e., if all the eigenvalues of the
matrix $D^2F_0(x_0)$ are strictly greater than $0$, then $x_0$ is a constrained
minimum for $f(x)$ with respect to the constraint $g(x) = 0$.

If the matrix $D^2F_0(x_0)$ is negative definite, i.e., if all the eigenvalues of the
matrix $D^2F_0(x_0)$ are strictly less than $0$, then $x_0$ is a constrained maximum
for $f(x)$ with respect to the constraint $g(x) = 0$.
\end{theorem}

Summarizing, the steps for solving the simplified version of a constrained
optimization problem using Lagrange multipliers are:

\medskip
\noindent\underline{Step 1:} Check that $\rank(Dg(x)) = m$ for all $x$ such that
$g(x) = 0$.

\medskip
\noindent\underline{Step 2:} Find $(x_0,\lambda_0) \in U \times \R^m$ such that
$D_{(x,\lambda)}F(x_0,\lambda_0) = 0$.

\medskip
\noindent\underline{Step 3:} Check whether the matrix $D^2F_0(x_0)$ is positive
semidefinite or negative semidefinite.

\medskip
\noindent\underline{Step 4:} Use Theorem 10.7 to decide whether $x_0$ is a
constrained minimum point or a constrained maximum point.

\medskip
The Lagrange multipliers method is used, for example, for portfolio optimization
problems; see sections 9.3, 9.4, and 9.5 for details.

\subsection{The ``Big O'' notation}

In this book, the ``Big O'' notation is primarily used for polynomials; a formal
definition can be found below. More details on the general form of the ``Big O''
notation can be found in Section 10.5 of Stefanica \cite{36}.

\begin{definition}
Let $f : \R \to \R$, and let $k$ be a positive integer. We write that
$f(n) = O(n^k)$, as $n \to \infty$, if and only if there exist constants
$C_1,C_2 > 0$ and $M > 0$ such that
\[
C_1 \;\le\; \left|\frac{f(n)}{n^k}\right| \;\le\; C_2, \qquad \forall\, n \ge M.
\]
\end{definition}

In particular, note that
\begin{equation}
\text{if} \quad 0 \;<\; \lim_{n\to\infty}\left|\frac{f(n)}{n^k}\right| \;<\; \infty,
\quad \text{then} \quad f(n) = O(n^k), \quad \text{as } n \to \infty.
\end{equation}

Let
\[
P(x) \;=\; \sum_{i=0}^p a_ix^i \;=\; a_px^p + a_{p-1}x^{p-1} + \ldots + a_1x + a_0,
\]
with $a_p \ne 0$, be a polynomial of degree $p$. Then, from (10.74), it follows that
\[
P(n) \;=\; O(n^p), \quad \text{as } n \to \infty.
\]

More precisely, if $a_{p-1} \ne 0$, then,
\begin{equation}
P(n) \;=\; a_px^p \;+\; O(n^{p-1}), \quad \text{as } n \to \infty.
\end{equation}

To prove, e.g., (10.75), note that
\begin{align*}
\lim_{n\to\infty}\left|\frac{P(n)-a_px^p}{n^{p-1}}\right|
&= \lim_{n\to\infty}\left|\frac{a_{p-1}n^{p-1}+a_{p-2}n^{p-2}+\ldots+a_1n+a_0}
{n^{p-1}}\right| \\
&= \lim_{n\to\infty}\left|a_{p-1} \;+\; \frac{a_{p-2}}{n} \;+\; \ldots \;+\;
\frac{a_1}{n^{p-2}} \;+\; \frac{a_0}{n^{p-1}}\right| \\
&= |a_{p-1}| \;>\; 0,
\end{align*}
and therefore $P(n) - a_px^p = O(n^{p-1})$; cf. (10.74)

The results below follow from (10.75) and are used in the book:
\begin{align}
8n - 7 &= 8n \;+\; O(1), \quad \text{as } n \to \infty; \\
10n - 7 &= 10n \;+\; O(1), \quad \text{as } n \to \infty; \\
n^2 + n - 1 &= n^2 \;+\; O(n), \quad \text{as } n \to \infty; \\
\frac{2n^3}{3} - \frac{n^2}{2} + \frac{5n}{6} - 1 &= \frac{2}{3}n^3 \;+\; O(n^2),
\quad \text{as } n \to \infty; \\
\frac{n^3}{3} + \frac{n^2}{2} + \frac{n}{6} &= \frac{1}{3}n^3 \;+\; O(n^2), \quad
\text{as } n \to \infty.
\end{align}

Rules for operations with the ``Big O'' notations can be derived from Definition
10.5 and are as follows:

If $k$ and $j$ are positive integers and $c \ne 0$ is a constant, then
\begin{align}
c\,n^jO(n^k) &= O(n^{j+k}), \quad \text{as } n \to \infty; \\
c\,O(n^k) &= O(n^k), \quad \text{as } n \to \infty; \\
O(n^k) + cn^k &= O(n^k), \quad \text{as } n \to \infty; \\
O(n^k) + O(n^k) &= O(n^k), \quad \text{as } n \to \infty; \\
O(n^k) + O(n^j) &= O(n^k), \quad \text{if } j < k, \quad \text{as } n \to \infty.
\end{align}

The following results can be proved using (10.81--10.84):

As $n \to \infty$,
\begin{align}
\left(\frac{2}{3}n^3+O(n^2)\right) + 2\left(n^2+O(n)\right)
&= \frac{2}{3}n^3+O(n^2)+2n^2+2O(n) \notag \\
&= \frac{2}{3}n^3+O(n^2)+O(n) \notag \\
&= \frac{2}{3}n^3 \;+\; O(n^2); \\
\left(\frac{1}{3}n^3+O(n^2)\right) + 2\left(n^2+O(n)\right)
&= \frac{1}{3}n^3 \;+\; O(n^2);
\end{align}
\begin{align}
\frac{2}{3}n^3+O(n^2) + n\left(2n^2+2O(n)\right)
&= \frac{8}{3}n^3+O(n^2)+2nO(n) \notag \\
&= \frac{8}{3}n^3+O(n^2)+O(n^2) \notag \\
&= \frac{8}{3}n^3 \;+\; O(n^2); \\
\frac{1}{3}n^3+O(n^2) + n\left(2n^2+2O(n)\right)
&= \frac{7}{3}n^3 \;+\; O(n^2).
\end{align}

\section{European options overview}

The material in this section is adapted from Stefanica \cite{36}; see sections
1.7--1.10 and 3.7 therein for more details.

A \textbf{Call Option} on an underlying asset (e.g., on one share of a stock, for an
equity option\footnote{The underlying asset for equity options is usually 100
shares, not one share. For clarity and simplicity reasons, we will be consistent
throughout the book in our assumption that options are written on just one unit of
the underlying asset.}) is a contract between two parties which gives the buyer of
the option the right, but not the obligation, to \textbf{buy} from the seller of the
option one unit of the asset (e.g., one share of the stock) at a predetermined time
$T$ in the future, called the \textbf{maturity} of the option, for a predetermined
price $K$, called the \textbf{strike} of the option. For this right, the buyer of
the option pays $C(t)$ at time $t < T$ to the seller of the option.

A \textbf{Put Option} on an underlying asset is a contract between two parties which
gives the buyer of the option the right, but not the obligation, to \textbf{sell} to
the seller of the option one unit of the asset at a predetermined time $T$ in the
future, called the \textbf{maturity} of the option, for a predetermined price $K$,
called the \textbf{strike} of the option. For this right, the buyer of the option
pays $P(t)$ at time $t < T$ to the seller of the option.

The options described above are plain vanilla European options. An option which can
be exercised at any time prior to maturity is called an American option.

Let $S(t)$ be the price of the underlying asset at time $t$. A call option is
at--the--money (ATM) if its strike is equal to the spot price, i.e., if $K = S(t)$.
Similarly, a put option is at--the--money (ATM) if its strike is equal to the spot
price, i.e., if $K = S(t)$. A call option is in--the--money (ITM) or
out--of--the--money (OTM) at time $t$ if $S(t) > K$ or $S(t) < K$, respectively. A
put option is in--the--money or out--of--the--money at time $t$ if $S(t) < K$ or
$S(t) > K$, respectively.

The payoff of a call option at maturity is
\begin{equation}
C(T) \;=\; \max(S(T)-K,0) \;=\; \begin{cases} S(T)-K, & \text{if} \;\; S(T) > K; \\
0, & \text{if} \;\; S(T) \le K. \end{cases}
\end{equation}

The payoff of a put option at maturity is
\begin{equation}
P(T) \;=\; \max(K-S(T),0) \;=\; \begin{cases} 0, & \text{if} \;\; S(T) \ge K; \\
K-S(T), & \text{if} \;\; S(T) < K. \end{cases}
\end{equation}

\medskip
\noindent\textbf{The Put--Call parity}
